\section{Aturan Konjungsi, Penambahan, dan Penyederhanaan}

Selain aturan yang berkaitan dengan implikasi, diperlukan pula aturan untuk operator konjungsi dan disjungsi. Tiga aturan berikut sering dipakai dalam pembuktian berantai: konjungsi, penambahan (\textit{addition}), dan penyederhanaan (\textit{simplification}).

\subsection{Aturan Konjungsi}

Jika dua proposisi masing-masing telah ditetapkan benar, keduanya dapat digabung menjadi satu konjungsi.

\begin{teorema}[Konjungsi]
Jika $p$ benar dan $q$ benar, maka $p \land q$ benar.
\end{teorema}

\textbf{Skema argumen Konjungsi:}
\[
\begin{array}{c}
p \\
q \\
\hline
\therefore p \wedge q
\end{array}
\]

\begin{contoh}
Premis 1: Jakarta adalah ibu kota Indonesia. ($p$) \\
Premis 2: Bangkok adalah ibu kota Thailand. ($q$) \\
Konklusi: Jakarta adalah ibu kota Indonesia, dan Bangkok adalah ibu kota Thailand. ($p \land q$)
\end{contoh}

\subsection{Aturan Penambahan Disjungsi (\textit{Addition})}

Aturan ini memungkinkan penambahan disjungsi terhadap proposisi yang telah ditetapkan benar.

\begin{teorema}[Penambahan Disjungtif]
Jika $p$ benar, maka $p \vee q$ benar untuk sembarang proposisi $q$.
\end{teorema}

\textbf{Skema argumen Penambahan:}
\[
\begin{array}{c}
p \\
\hline
\therefore p \vee q
\end{array}
\]

\begin{contoh}
Premis: ``Budi berolahraga setiap pagi.'' ($p$) \\
Konklusi: ``Budi berolahraga setiap pagi, atau hujan turun di kutub selatan hari ini.'' ($p \vee q$) \\
Konklusi tetap valid karena $p$ benar menjamin $p \vee q$ benar, terlepas dari nilai kebenaran $q$.
\end{contoh}

\subsection{Aturan Penyederhanaan Konjungsi (\textit{Simplification})}

Dari konjungsi yang benar, salah satu komponennya dapat diambil sebagai konklusi.

\begin{teorema}[Penyederhanaan]
Jika $p \land q$ benar, maka $p$ benar dan $q$ benar. Oleh karena itu, dapat disimpulkan $p$ saja atau $q$ saja.
\end{teorema}

\textbf{Skema argumen Penyederhanaan:}
\[
\begin{array}{ccc}
p \wedge q & \text{atau} & p \wedge q\\
\hline
\therefore p & & \therefore q
\end{array}
\]

\begin{contoh}
Premis: ``Air mendidih pada suhu 100$^\circ$C pada tekanan standar, dan api dapat membakar bahan mudah terbakar.'' ($p \land q$) \\
Konklusi yang valid: ``Air mendidih pada suhu 100$^\circ$C pada tekanan standar.'' ($\therefore p$), atau ``Api dapat membakar bahan mudah terbakar.'' ($\therefore q$) \cite{rosen2019}.
\end{contoh}
